\section{Definition of the Study Context}
\subsection{Country Selection for Comparative Analysis}
This study evaluates four European countries: France, Italy, Poland, and Switzerland, to assess how national energy contexts influence the ranking of nuclear reactor designs \footnote{Data regarding electricity loads and cross border flows is source from the ENTSO-E Transparency Platform and refers to data from 2024 \cite{ENTSOE_Transparency_2025}. Data analysis can be found at \textcolor{red}{ADD LINK !!!!!!!!}}.
--- try to justify why Europe and why select some... Still, there is no enough rationale for selecting these four countries---

\textbf{Italy}'s nuclear landscape is constrained by the 1987 post-Chernobyl referendum which mandated complete phase-out of nuclear power and its 2011 reaffirmation following Fukushima. However, since Italian referendums are abrogative rather than prescriptive, government decrees are enough to reverse these policies \cite{Costituzione_Italiana_Art75_1947}. The current administration seems willing to reintroduce nuclear power \cite{Sole24Ore_Italy_Nuclear_Strategy_2024}, however, Italy's history of political volatility and the long-term nature of nuclear projects create significant implementation challenges regardless of public opinion trends. 

For this study we have taken into consideration only the northern part of Italy due to the country's fragmented electricity market structure. The north bidding zone is the most energy-intensive with an average load of 17.66 $\pm$ 4.35 GW, for reference, by analysing data from the Italian \gls{tso} \cite{Terna_Italian_Electricity_Demand_2024} and the Italian institute for statistics \cite{ISTAT_POS_2024}, it is possible to conclude that the North bidding zone, despite only hosting 35\% of the Italian population, represents 56\% of the total electric demand. In addition to the significant load, the bidding zone has a substantial dependence on imports, which averaged 5.23 $\pm$ 1.67 GW in 2024, accounting for almost 30\% of its electric demand. Analysis of cross-border flow data shows that imports are active 99\% hours of the year. 

\paragraph{Poland} shows strong public backing for nuclear energy development, with surveys reporting up to $92\%$ support \cite{NucNet_Poland_Nuclear_Support_2024}. The public sentiment aligns with tangible progress in the country's nuclear program, including preliminary site works and vendor agreements \cite{Westinghouse_Poland_Nuclear_2025,WorldNuclearNews_Poland_Preliminary_Works_2025}, suggesting robust governmental commitment to nuclear deployment. The Polish grid produced 55\% of its electricity with coal \cite{IEA_Electricity_2024}, making it the country with the highest green house gas emission intensity in Europe \cite{ElectricityMaps_CarbonIntensity_2024}.


\paragraph{Switzerland} presents a more complex picture. Following the 2011 Fukushima accident, the government adopted a phase-out policy later confirmed by a referendum in 2017 \cite{WorldNuclear_Switzerland_Profile_2025}. However, recent surveys indicate shifting public opinion, with $53\%$ now supporting nuclear power \cite{Euronuclear_Swiss_Nuclear_Poll_2023,NucNet_Nuclear_Support_Survey_2023,WorldNuclearNews_Swiss_Nuclear_Poll_2023}. The federal government has subsequently signalled potential reconsideration of nuclear role in meeting climate targets \cite{swissinfo_Nuclear_Power_Switzerland_2025}, though the narrow margin of support and Switzerland's direct democracy system create policy uncertainty. 

The Swiss grid is characterized by a low demand, averaging 6.79 $\pm$ 1.12 GW, and generation dominated by renewables (65\% average share) particularly hydroelectric which provided 59\% of the electricity in 2024 \cite{IEA_Electricity_2024}. This high dependence on renewables is reflected on substantial cross border flow variability which average 2.71 $\pm$ 3.14 GW. The large standard deviation reflects the oscillating nature of the Swiss cross-border flows, consequence of periodic exports due to overproduction, alternating with imports when renewable output is low.\\

\paragraph{France} demonstrates consistent public support for nuclear energy, with 75\% favouring continued operation of existing plants and $65\%$ supporting new construction according to a 2022 survey \cite{IFOP_Nuclear_Perception_2023}. Given the long-standing consensus on nuclear power as a strategic energy pillar and the country's established industrial base, significant policy reversals appear unlikely. The country electric grid relies heavily on nuclear power, which accounted for 67\% of its electricity generation \cite{IEA_Electricity_2024}. Despite the substantial demand, which averaged $49.05 \pm 9.83$ GW, France is a net electricity exporter, with 21\% of total domestic production crossing the border to neighbouring countries.

\subsection{Selection of Reactor Models}
This study focuses on a curated set of reactor designs at an advanced stage of development.
The selection process begins with the analysis of the \gls{iaea} "Nuclear Reactors in the World" publication \cite{international2024iaea}, which catalogues operational, under-construction, and decommissioned reactors globally. \\

The initial selection process applied two primary exclusion criteria. First, designs were eliminated based on supplier origin, removing all reactors associated with countries unlikely to engage in nuclear technology transfer with European countries, namely, Russia, China and India. Second, reactors commissioned before 2000 were excluded from consideration, as their designs represent outdated technological standards \footnote{The complete filtering is documented in the accompanying analysis notebook \textcolor{red}{ADD LINK !!!!!!!!}}.

From this pre-selected group, three designs were prioritized for comprehensive analysis: the \gls{epr}-1750, \gls{apr}, and \gls{ap}. It should be noted that while the \gls{epr} (1650 MWe) and \gls{epr}-1750 share significant design similarities, they are treated as separate models in this study. Additionally, the \gls{abwr}, though listed as operational in the \gls{iaea} \gls{aris} database, was excluded following verification through the \gls{pris} database, which indicates that all \gls{abwr} units have suspended operation since 2011 \cite{IAEA_PRIS_2025}. The selection process and its outcomes are summarized in Table \ref{tab:reactor-selection}.

The author acknowledges the controversy between Westinghouse and KEPCO regarding the APR design derivation from the System 80 design. However, as no official documentation prohibits its commercialization in Europe, and considering recent evaluations of the Korean design (including the 2023 Czech tender, though later contested by \gls{edf} and Westinghouse \cite{WNN_Czech_Nuclear_Appeal_2024}), the \gls{apr} has been included in this analysis.

The \gls{epr}2 and \gls{epr}1200 designs, while reviewed positively by the French nuclear safety authority \cite{WNN_EPR1200_Safety_2023}, lack commercial commitments (the gls{epr}1200 was proposed for the Czech tender but ultimately unsuccessful \cite{WNN_Czech_Nuclear_Appeal_2024}). Moreover, considering their absence from the ARIS database and overall limited available technical documentation, these \gls{edf} designs have been excluded from the current study.


\begin{table}[!hb]
\centering
\begin{tabularx}{\textwidth}{llXXXX}
\toprule
\textbf{Reactor Name} & \textbf{Country} & \textbf{Technology-sharing} & \textbf{Commercial Operation} & \textbf{Up to Date} & \textbf{Decision} \\
\midrule
\gls{apr} & S.Korea & \cmark & \cmark & \cmark & \textbf{Selected} \\

CAREM Prototyp & Argentina & \cmark & \xmark & \cmark & Rejected \\

PRE KONVOI & Germany & \cmark & \cmark & \xmark & Rejected \\

\gls{epr} 1750 & France & \cmark & \cmark & \cmark & \textbf{Selected} \\

\gls{ap} & USA & \cmark & \cmark & \cmark & \textbf{Selected} \\

\gls{abwr} & Japan & \cmark & \xmark & \cmark & Rejected \\

\gls{apwr} & Japan & \cmark & \xmark & \cmark & Rejected \\

\gls{epr} & France & \cmark & \cmark & \xmark & Rejected \\

OPR-1000 & S.Korea & \cmark & \cmark & \xmark & Rejected \\
\bottomrule
\end{tabularx}
\caption{Selection of reactor models for analysis}
\label{tab:reactor-selection}
\end{table}


Subsequently, the \gls{iaea} \gls{aris} database \cite{iaea2025aris}, was consulted to identify advanced reactor models classified as under design, under regulatory review, under construction, or operational. To avoid redundancy, Generation III+ designs captured in both \gls{aris} and \gls{pris} datasets were removed. Demonstrations units and prototypes were also excluded from this set, yielding 10 candidate designs. Each of them was then assessed for development status, technological readiness, and applicability to European deployment as summarized in Table \ref{tab:advanced-reactor-selection}.

\begin{table}[!hb]
\centering
\begin{tabularx}{\textwidth}{lllX}
\toprule
\textbf{Reactor Name} & \textbf{Country} & \textbf{Used} & \textbf{Reasoning} \\
\midrule

APR+ & S.Korea & \xmark & As reported in the manufacturer website \cite{kepcoenc2025} this design is mean to be an "indigenous reactor cleared of obstacles to export"\\

ATMEA & Japan, France & \xmark & Joint venture was abandoned in 2019 \cite{PLACEHOLDER}\\

BWRX300 & USA, Canada & \cmark & Design from GE-Hitachi, a lot of support from US and Canada as well as interes in the EU \cite{wnn2025bwrx300} \cite{wnn2025hungarysmr}.\\

i-SMR & S.Korea & \xmark & Never expressed interest in development outside south Asia and the middle east, no updates since early 2024.\\

iMSR400 & US & \xmark & No interest overseas, project seems to be at an early stage of development.\\

Natrium & US & \xmark & Demo plant \cite{terrapower2025natrium}.\\

Nuward & France & \xmark & After the change in direction of 2024 the project has been set back \cite{PLACEHOLDER}.\\

Nuscale & US & \cmark & Even if there was some financial controversy in 2024 \cite{PLACEHOLDER}. the project seems to be still moving forward in terms of goverment founding \cite{PLACEHOLDER} and regulatory review \cite{PLACEHOLDER}.\\

PRISM & US & \xmark & GE-Hitachi considers it a "concept ... currently being put into practice in two reactors: Natrium and ARC-100" \cite{PLACEHOLDER}.\\

Rolls Royce SMR & UK & \cmark & It's moving forward and from the technological point of view this is the simples as it it litteraly a scaled down PWR\cite{PLACEHOLDER}.\\

SMART & S.Korea, Saudi Arabia & \xmark & Plan to export only to south-east asia and middle east \cite{wnn2025smrdesign}.\\ \bottomrule

\end{tabularx}
\caption{Selection of advanced reactor models for analysis}
\label{tab:advanced-reactor-selection}
\end{table}

\subsection{Overview of the case study}
The overall selection produced a set of six reactors, three Generation III+ and three \glspl{smr}. Their main characteristics are summarized in Table \ref{tab:final-reactor-selection}. The most notable differences are the integrated designs of the \gls{bwrx} and Nuscale \gls{smr} which eliminate external primary loops in favour of steam generators positioned inside the reactor pressure vessel. This allows to reduce the overall size and to exploit natural circulation for cooling. The \gls{npm}, referred to as US460 in technical documentation, is the only model designed to be deployed in a multi unit configuration, with up to 12 modules in the same containment building, for this study the 6 module plant (VOYGR-6) has been selected as it recently revived the \gls{sda} by the US Nuclear Regulatory Commission in May 2025 \footnote{\textit{"A standard design approval indicates that a proposed reactor design meets applicable agency safety requirements. Companies that seek to use the US460 design would have to file applications seeking permission to build and operate a nuclear reactor using the approved design"} \cite{NRC_2025_033}}.

\begin{table}[!hb]
\centering
\begin{tabularx}{\textwidth}{lclcX}
\toprule
\textbf{Reactor Name} & \textbf{Type} & \textbf{Construction} & \textbf{P$_e$ (MWe)} & \textbf{Supplier} \\
\midrule
\gls{apr} & \gls{pwr}  & 1 Unit & 1340 & KEPCO (S.Korea) \\
\gls{epr} 1750 & \gls{pwr}  & 1 Unit & 1630 & Framatome (France) \\
\gls{ap} & \gls{pwr}  & 1 Unit & 1120 & Westinghouse (USA) \\
\gls{bwrx} & \gls{ibwr} & 1 Unit & 300 & GE-Hitachi (USA)\\
Rolls Royce SMR & \gls{pwr}  & 4 to 12 Units & 470 & Rolls Royce (UK) \\
Nuscale & \gls{ipwr} & 1 Unit & 77 & Nuscale Power (USA) \\
\bottomrule
\end{tabularx}
\caption{Final selection of reactor models for analysis}
\label{tab:final-reactor-selection}
\end{table}

